\documentclass[12pt]{article}

% Core math
\usepackage{amsmath,amssymb,amsfonts}
\usepackage{amsthm}

% Diagrams
\usepackage{tikz-cd}
\usepackage{tikz}
\usetikzlibrary{calc}

% References
\usepackage{hyperref}
\usepackage{cleveref}

% Graphics
\usepackage{graphicx}

% Page layout
\usepackage[margin=1in]{geometry}

% GrokRxiv sidebar
\usepackage{xcolor}

% ==== topological-phases-of-matter : canonical macros (v1, 2026-07-14) ====
% Requires amsmath, amssymb, amsfonts (mathfrak, mathbf, mathcal all standard).

% --- condensation / probes ---
\newcommand{\cond}[1]{\underline{#1}}          % condensation X |-> \cond{X}
\newcommand{\cg}[1]{\underline{#1}}            % condensed group (alias, for readability)
\newcommand{\Cont}{\operatorname{Cont}}        % Cont(S,X)
\newcommand{\Shape}{\operatorname{Shape}}      % shape / homotopy type

% --- interaction space & locality ---
\newcommand{\Int}{\mathcal{I}}                 % \Int_{d,G} interaction space
\newcommand{\Ffun}{F}                          % F-function (interaction decay)
\newcommand{\BF}{\mathcal{B}_{F}}              % interaction Banach space for F

% --- moduli stacks ---
\newcommand{\Ham}{\mathfrak{Ham}}              % \Ham_{d,G}
\newcommand{\Gap}{\mathfrak{Gap}}              % \Gap_{d,G}
\newcommand{\Phase}{\mathfrak{Phase}}          % \Phase_{d,G}
\newcommand{\Phases}{\operatorname{Phases}}    % \Phases_{d,G} = pi_0 Shape Phase
\newcommand{\Weq}{\mathcal{W}}                 % class of equivalences to invert
\newcommand{\st}{\mathrm{st}}                  % stabilization superscript
\newcommand{\hfix}[1]{h#1}                     % homotopy-fixed-point superscript, e.g. \Phase_d^{\hfix{\cg{G}}}
\newcommand{\quotstack}[2]{[#1/#2]}            % action stack [Y/\cg{G}]

% --- gap ---
\newcommand{\gapof}[1]{\operatorname{gap}(#1)} % gap(H)
\newcommand{\gapbound}{\Delta}                  % uniform gap bound

% --- stacking / spectrum ---
% \boxtimes is built in (stacking product)
\newcommand{\IPcond}{\mathbf{IP}^{\mathrm{cond}}}  % invertible condensed phase spectrum

% --- transitions / discriminant / charge ---
\newcommand{\Sig}{\Sigma}                       % \Sig_f gapless discriminant
\newcommand{\rc}{\partial\nu}                   % relative transition charge in E^{q+1}(B,U_f)

% --- observables & solid K-theory ---
\newcommand{\Kop}{K_{\mathrm{op}}}             % operator / topological K-theory
\newcommand{\Kalg}{K_{\mathrm{alg}}}           % algebraic K-theory
\newcommand{\Solid}{\operatorname{Solid}}      % solidification functor

% --- disorder ---
\newcommand{\dis}{\Omega}                       % disorder hull \Omega = Q^{Z^d} (Q = local-configuration alphabet; NOT the F-function)
% ==== end canonical macros ====

% --- local convenience macros (not overriding canonical glyphs) ---
\newcommand{\R}{\mathbb{R}}
\newcommand{\C}{\mathbb{C}}
\newcommand{\Z}{\mathbb{Z}}
\newcommand{\N}{\mathbb{N}}
\newcommand{\Hilb}{\mathcal{H}}
\newcommand{\Alg}{\mathcal{A}}
\newcommand{\gsp}{\mathcal{G}}                 % ground-state space
\newcommand{\norm}[1]{\lVert #1 \rVert}
\newcommand{\abs}[1]{\lvert #1 \rvert}
\newcommand{\ip}[1]{\langle #1 \rangle}
\newcommand{\set}[1]{\{\,#1\,\}}
\newcommand{\gapfin}[1]{\operatorname{gap}_{#1}}
\newcommand{\Tr}{\operatorname{Tr}}
\newcommand{\supp}{\operatorname{supp}}
\newcommand{\dist}{\operatorname{dist}}
\DeclareMathOperator*{\einf}{inf}

% Theorem environments. We use aliascnt so that each environment carries its own
% counter (aliased to the Theorem counter for shared numbering); this lets cleveref
% name cross-references correctly ("Proposition 2.3", "Lemma 4.1") instead of
% collapsing them all to "Theorem".
\usepackage{aliascnt}
\theoremstyle{plain}
\newtheorem{theorem}{Theorem}[section]
\newaliascnt{proposition}{theorem}
\newtheorem{proposition}[proposition]{Proposition}
\aliascntresetthe{proposition}
\newaliascnt{lemma}{theorem}
\newtheorem{lemma}[lemma]{Lemma}
\aliascntresetthe{lemma}
\newaliascnt{corollary}{theorem}
\newtheorem{corollary}[corollary]{Corollary}
\aliascntresetthe{corollary}
% Named (unnumbered) principal results, tagged per the series convention.
\newtheorem*{thmA}{Theorem III-A}
\newtheorem*{thmB}{Theorem III-B}
\newtheorem*{thmC}{Theorem III-C}
\newtheorem*{conjone}{Conjecture III-1}
\newtheorem*{conjtwo}{Conjecture III-2}
\newtheorem*{conjthree}{Conjecture III-3}
\theoremstyle{definition}
\newaliascnt{definition}{theorem}
\newtheorem{definition}[definition]{Definition}
\aliascntresetthe{definition}
\newaliascnt{example}{theorem}
\newtheorem{example}[example]{Example}
\aliascntresetthe{example}
\theoremstyle{remark}
\newaliascnt{remark}{theorem}
\newtheorem{remark}[remark]{Remark}
\aliascntresetthe{remark}
% cleveref display names for the aliased environments.
\crefname{theorem}{Theorem}{Theorems}
\Crefname{theorem}{Theorem}{Theorems}
\crefname{proposition}{Proposition}{Propositions}
\Crefname{proposition}{Proposition}{Propositions}
\crefname{lemma}{Lemma}{Lemmas}
\Crefname{lemma}{Lemma}{Lemmas}
\crefname{corollary}{Corollary}{Corollaries}
\Crefname{corollary}{Corollary}{Corollaries}
\crefname{definition}{Definition}{Definitions}
\Crefname{definition}{Definition}{Definitions}
\crefname{example}{Example}{Examples}
\Crefname{example}{Example}{Examples}
\crefname{remark}{Remark}{Remarks}
\Crefname{remark}{Remark}{Remarks}

% GrokRxiv DOI sidebar
\definecolor{grokgray}{RGB}{110,110,110}
\AddToHook{shipout/foreground}{%
  \ifnum\value{page}=1
    \begin{tikzpicture}[remember picture, overlay]
      \node[
        rotate=90,
        anchor=south,
        font=\Large\sffamily\bfseries\color{grokgray},
        inner sep=0pt
      ] at ([xshift=38pt, yshift=0.52\paperheight]current page.south west)
      {GrokRxiv:2026.07.spectral-gap-stability\quad
       [\,math-ph\,]\quad
       14 Jul 2026};
    \end{tikzpicture}
  \fi
}

\hypersetup{
  colorlinks=true,
  linkcolor=blue!55!black,
  citecolor=green!45!black,
  urlcolor=blue!55!black
}

\title{\textbf{The Uniformly Gapped Substack:\\
Existence and Stability of the Thermodynamic Spectral Gap}\\[0.4em]
\large Part III of \emph{Topological Phases of Matter in the Condensed-Mathematics Paradigm}}

\author{Matthew Long \\
\textit{The YonedaAI Collaboration} \\
\textit{YonedaAI Research Collective} \\
Chicago, IL \\
\texttt{matthew@yonedaai.com} $\cdot$ \url{https://yonedaai.com}}

\date{14 July 2026}

\begin{document}
\maketitle

\begin{abstract}
Fix a spatial dimension $d$, a lattice $L$, local Hilbert spaces, a locality class,
and a symmetry group $G$. The condensed moduli stack of $G$-symmetric quasi-local
Hamiltonians $\Ham_{d,G}$ carries a distinguished substack $\Gap_{d,G}$, the
\emph{uniformly} gapped systems: over a profinite probe $S$, a family lies in
$\Gap_{d,G}(S)$ exactly when $\einf_{s\in S}\gapof{H_s}\ge\gapbound$ for some
$\gapbound>0$. We take the word \emph{uniformly} seriously (a family whose
pointwise gaps degenerate to zero is not a member) and study $\Gap_{d,G}$ as the
object on which the phase functor of the program rests. The governing constraint is
a hard theorem: by Cubitt, P\'erez-Garc\'ia and Wolf the spectral gap is
undecidable, so there is no algorithm and no uniform criterion that decides
membership in $\Gap_{d,G}$ for general translation-invariant families. We therefore
do not attempt to \emph{decide} the gap. Existence of a thermodynamic gap is treated
as a hypothesis that \emph{defines} the substack, and the content of the paper is
\textbf{stability and structure}, not decidability. We prove three results from
cited present-day technology. First, at fixed finite volume the $n$-gap is
Lipschitz in the interaction norm (with constant $2C_\Lambda$), so the gapped locus is open at every finite
resolution and the gap is continuous along continuous families (Theorem III-A)\@.
Second, on the frustration-free / LTQO stratum the Bravyi--Hastings--Michalakis and
Michalakis--Zwolak stability theorems, with the Nachtergaele--Sims--Young
infinite-volume bulk-gap result, supply a volume-independent perturbation threshold,
which we recast as: $\Gap_{d,G}$ contains a condensed-open neighborhood of each
stratum point along quasi-local perturbation directions (Theorem III-B)\@. Third, a
uniform gap forces uniform exponential clustering, with a correlation length bounded
uniformly over the profinite base (Theorem III-C)\@. Beyond the stratum we record what
we cannot prove as three numbered conjectures: openness of $\Gap_{d,G}$ in the full
condensed topology, detected by finite quotients (III-1); descent of uniform
gappedness along profinite disorder hulls $\dis=Q^{\Z^d}$ (III-2); and the
identification of quasi-adiabatic continuation with the internal path components
realizing $\Weq$, so that $\pi_0\Shape\!\big(\Gap_{d,G}[\Weq^{-1}]\big)$ equals the
operator-algebraic set of gapped ground-state phases (III-3). A transverse-field
Ising computation renders the gapless discriminant $\Sig=\{g=1\}$ as a numerical
picture and exhibits, away from criticality, both gap persistence and the Lipschitz
bound of Theorem III-A.
\end{abstract}

\tableofcontents

\section{Introduction}
\label{sec:intro}

\subsection{The gapped substack and why uniformity is the point}

The program of which this is the third part organizes topological phases of matter
around a single geometric object: the condensed moduli stack $\Ham_{d,G}$ of
$G$-symmetric local and quasi-local Hamiltonians on a lattice $L$ in spatial
dimension $d$. Part I \cite{paperLocality} makes the interaction space into a Banach
space $\BF$ through an $\Ffun$-function locality norm and promotes the Heisenberg
dynamics to a morphism of condensed objects; Part II \cite{paperPositivity} attaches
the quasi-local $C^*$-algebra, its compact state space, and the solid $K$-theory
invariant of \cite{aoki-solidification}. The present paper concerns the substack

\begin{equation}
\label{eq:gap-def-intro}
\Gap_{d,G}(S)\;=\;\bigcup_{\gapbound>0}\;
\set{H\in\Ham_{d,G}(S)\,:\,\einf_{s\in S}\gapof{H_s}\ge\gapbound},
\qquad S\ \text{profinite},
\end{equation}

the \emph{uniformly} gapped systems.\par\noindent
Everything downstream (the stabilized phase
$\infty$-groupoid $\Phase_{d,G}=\big(\Gap_{d,G}[\Weq^{-1}]\big)^{\st}$, the phase
set $\Phases_{d,G}=\pi_0\Shape(\Phase_{d,G})$, and the invertible condensed phase
spectrum $\IPcond_{d,G}$ built in Part IV \cite{paperEFT}) is a construction on top
of \eqref{eq:gap-def-intro}. The phase functor is only as good as the substack it is
built from.

The word \emph{uniformly} is not decorative. A profinite family $s\mapsto H_s$ in
which each $H_s$ happens to be gapped, but with $\einf_{s}\gapof{H_s}=0$, is exactly
the kind of family that fails to define a robust invariant: as one moves through the
base the gap can be driven arbitrarily small, correlations grow without bound, and
the locally constant phase label $\nu_f$ of the program loses meaning at the
degenerating points. The union over $\gapbound>0$ in \eqref{eq:gap-def-intro} records
that a member of $\Gap_{d,G}(S)$ is a family for which \emph{some} strictly positive
$\gapbound$ bounds the gap below \emph{simultaneously} over the whole probe. Keeping
the $\einf_{s}$ explicit is the single most important discipline of this paper.

\subsection{The undecidability wall}
\label{sec:intro-wall}

There is a hard theorem standing between us and any naive program to ``compute the
substack.'' Cubitt, P\'erez-Garc\'ia and Wolf \cite{cpw-undecidable} proved that the
spectral gap is \emph{undecidable}: there is a family of translation-invariant,
nearest-neighbour Hamiltonians on a two-dimensional lattice, depending
computably on a parameter, for which no algorithm decides whether the
thermodynamic-limit system is gapped or gapless. The construction embeds the halting
problem into the low-energy physics, so the question ``is this system gapped?'' is at
least as hard as the halting problem, and in fact $\Pi_1$-complete in the arithmetic
hierarchy. The undecidability survives many natural restrictions.

We state the consequence for this paper in the strongest terms, because it shapes
every claim we make.

\begin{quote}
\emph{There is no general algorithm, decision procedure, or uniform criterion that
determines membership in $\Gap_{d,G}$ for arbitrary members of $\Ham_{d,G}$.
Existence of a thermodynamic gap is not, in general, a decidable property; it is an
\textbf{assumption that defines} the substack $\Gap_{d,G}$.}
\end{quote}

This is why the paper is titled around \emph{stability} rather than existence. We do
not offer a test for gappedness: no such test can exist in general. Instead we take
gappedness as a standing hypothesis and ask a question undecidability does not
forbid: \emph{given} a uniformly gapped family, what can be said about its
neighborhood, its correlations, and its behavior along the profinite base? The
answers are stability theorems, and they are the mathematically load-bearing content
of Part III\@. Any statement in this paper that could be read as deciding, bounding
from below, or algorithmically certifying the gap of a general member of $\Ham_{d,G}$
would contradict \cite{cpw-undecidable} and is, by construction, absent.

\subsection{Contributions}
\label{sec:intro-contrib}

Within the boundary set by \Cref{sec:intro-wall}, we prove three theorems and
formulate three conjectures.

\begin{itemize}
\item \textbf{Theorem III-A} (\Cref{sec:finite}). At a fixed finite volume the
$n$-gap of a Hamiltonian is $2C_\Lambda$-Lipschitz in the interaction norm (equivalently $2$-Lipschitz in the operator norm); consequently the
gapped locus is open at every finite resolution, and along any $S$-continuous family
the map $s\mapsto\gapof{H_s^{\Lambda}}$ is continuous. This is elementary (Weyl's
perturbation inequalities plus finite-volume norm equivalence), and it is the one
place where openness of $\Gap_{d,G}$ is unconditional. It is also the finite-volume
statement that undecidability leaves untouched: a finite matrix has a computable gap;
undecidability is a statement about the thermodynamic limit.

\item \textbf{Theorem III-B} (\Cref{sec:stability}). On the frustration-free / LTQO
stratum, the stability theorems of Bravyi--Hastings--Michalakis
\cite{bhm-stability,bravyi-hastings-shortproof}, Michalakis--Zwolak
\cite{michalakis-zwolak}, and Nachtergaele--Sims--Young \cite{nsy-bulkgap} provide a
perturbation threshold $\varepsilon_0>0$ independent of system size. We recast this
as a statement about $\Gap_{d,G}$: each point of the stratum has a condensed-open
neighborhood, of interaction-norm radius $\varepsilon_0$ along quasi-local
perturbation directions, lying inside $\Gap_{d,G}$. Openness and stability of the
substack hold \emph{on the stratum}; we are careful to claim nothing beyond it.

\item \textbf{Theorem III-C} (\Cref{sec:clustering}). A uniform gap over a family
forces uniform exponential clustering. Combining Hastings--Koma \cite{hastings-koma}
and Nachtergaele--Sims \cite{nachtergaele-sims-clustering} with the uniform
Lieb--Robinson velocity of Part I, a family with $\einf_s\gapof{H_s}\ge\gapbound$ has
a correlation length $\xi$ bounded above uniformly over the profinite base. This is
the structural feature separating uniformly gapped families from merely
pointwise-gapped ones.

\item \textbf{Conjectures III-1, III-2, III-3} (\Cref{sec:conjectures}). Openness of
$\Gap_{d,G}$ in the full condensed topology on the physically relevant stratum, with
the uniform-gap sheaf condition detected by finite quotients (III-1); descent of
uniform gappedness along profinite covers, i.e.\ disorder hulls $\dis=Q^{\Z^d}$
(III-2); and the identification of quasi-adiabatic continuation with the internal
path components realizing $\Weq$, so that
$\pi_0\Shape\!\big(\Gap_{d,G}[\Weq^{-1}]\big)$ equals the operator-algebraic set of
gapped phases of Ogata \cite{ogata-classification-review} (III-3).
\end{itemize}

\Cref{sec:numerics} illustrates the whole picture on the transverse-field Ising
chain: exact diagonalization renders the gapless discriminant $\Sig=\{g=1\}$, the
gap persists away from it in the manner Theorem III-A predicts, and, crucially, no
finite computation decides the thermodynamic gap, exactly as
\Cref{sec:intro-wall} requires.

\subsection{Relation to companion papers}
\label{sec:intro-companions}

This is Part III of a six-part series developing topological phases of matter in the
condensed-mathematics paradigm. The parts are modular: each takes the previous as
input and produces new structure.

\emph{Part I}, \emph{Condensed Locality: Lieb--Robinson Estimates and Quasi-Local
Dynamics on the Moduli Stack of Hamiltonians} \cite{paperLocality}, supplies the
$\Ffun$-function Banach space $\BF$ and the uniform Lieb--Robinson velocity that
Theorem III-C consumes; the quasi-adiabatic and automorphic-equivalence tools it
formalizes are the ones Conjecture III-3 invokes.

\emph{Part II}, \emph{Positivity, $C^*$-Norms, and Condensed State Spaces of
Quasi-Local Algebras} \cite{paperPositivity}, provides the infinite-volume GNS
framework in which the thermodynamic gap of \Cref{sec:framework} is defined, and the
crossed-product observable algebras through which the disorder hull of Conjecture
III-2 acquires its $K$-theory.

\emph{Part IV}, \emph{From Lattice Models to Effective Field Theories: Stabilization
and the Invertible Condensed Phase Spectrum} \cite{paperEFT}, is the immediate
consumer of this paper: it group-completes the invertible sector of
$\Phase_{d,G}=(\Gap_{d,G}[\Weq^{-1}])^{\st}$ into $\IPcond_{d,G}$. Its constructions
presuppose that $\Gap_{d,G}$ is a well-behaved, stable substack, which is what
Theorems III-A--III-C, on their stratum, secure.

\emph{Part V}, \emph{Physical Realizability of Bordism and Homotopy Classes by Gapped
Lattice Systems} \cite{paperRealizability}, asks which abstract classes are realized
by uniformly gapped lattice families; the ``uniformly gapped'' hypothesis it uses is
the one defined here.

\emph{Part VI}, \emph{Topological Phases of Matter in the Condensed-Mathematics
Paradigm: A Modular Research Program} \cite{paperSynthesis}, assembles the five
modules and states the global Master Conjecture, of which Conjectures III-1--III-3
are the gap-theoretic components.

\section{Mathematical framework}
\label{sec:framework}

We recall only what is needed and fix notation consistent with the series. Details of
the condensed formalism are in \cite{clausen-scholze-condensed,barwick-haine-pyknotic}
and Part I; the operator-algebraic framework is in
\cite{ogata-classification-review} and Part II\@.

\subsection{Condensed probes and the Hamiltonian stack}
\label{sec:framework-condensed}

A \emph{condensed set} is a sheaf on the site of profinite sets with finite,
jointly-surjective covers \cite{clausen-scholze-condensed}. For a topological space
$X$ its condensation $\cond{X}$ is defined on a profinite probe $S$ by
$\cond{X}(S)=\Cont(S,X)$; on compactly generated spaces $X\mapsto\cond{X}$ is fully
faithful, so ordinary parameter spaces, tori, Banach spaces such as $\BF$, and
compact disorder hulls all embed into a category with good algebraic and homological
behavior. We use \emph{light} condensed sets throughout to avoid set-theoretic size
issues in the profinite inverse limits \cite{clausen-scholze-analytic-stacks}.

For a profinite probe $S$, an element of $\Ham_{d,G}(S)$ is an $S$-continuous family
$s\mapsto\Phi_s$ of $G$-symmetric quasi-local interactions with finite
$\Ffun$-norm: equivalently, by Part I, a continuous map $S\to\BF$ landing in the
$G$-symmetric interactions. Each interaction $\Phi$ assigns to a finite $X\subset L$
a self-adjoint $\Phi(X)$ supported on $X$, and generates finite-volume Hamiltonians
$H_\Lambda(\Phi)=\sum_{X\subseteq\Lambda}\Phi(X)$ for finite $\Lambda\subset L$ and a
thermodynamic-limit dynamics and ground-state theory in the GNS sense of Part II\@. We
write $H$ for the interaction, $H_s$ for the member at $s\in S$, and $H_\Lambda$ for
a finite-volume restriction when the interaction is fixed.

\subsection{Finite-volume gap and thermodynamic gap}
\label{sec:framework-gap}

Two notions of gap appear, and their relationship is where the subtlety lives.

At \emph{finite volume} $\Lambda$, $H_\Lambda$ is a self-adjoint operator on the
finite-dimensional Hilbert space $\Hilb_\Lambda=\bigotimes_{x\in\Lambda}\Hilb_x$. Let
its eigenvalues, counted with multiplicity and in non-decreasing order, be
$\lambda_0(H_\Lambda)\le\lambda_1(H_\Lambda)\le\cdots$. For an integer $n\ge 1$ we
define the \emph{$n$-gap}
\begin{equation}
\label{eq:ngap}
\gapfin{\Lambda,n}(H)\;=\;\lambda_n(H_\Lambda)-\lambda_0(H_\Lambda),
\end{equation}
the spectral gap above the lowest $n$ levels. When the ground state is non-degenerate
one takes $n=1$; when there is a ground-state \emph{space} of dimension $m$ (as in
topological order or symmetry breaking) the physically relevant gap is the
$m$-gap, the distance from the ground-state sector to the first genuinely excited
level. The device of a fixed integer $n$ sidesteps the discontinuity that a changing
ground-state multiplicity would otherwise introduce; we return to this in
\Cref{sec:finite}.

At the \emph{thermodynamic} level, following Part II and
\cite{ogata-classification-review,nsy-bulkgap}, one works with an infinite-volume
ground state $\omega$, its GNS representation $(\Hilb_\omega,\pi_\omega,\Omega_\omega)$,
and the GNS Hamiltonian $H_\omega\ge 0$ implementing the dynamics with
$H_\omega\Omega_\omega=0$. The \emph{thermodynamic gap} is
\begin{equation}
\label{eq:thermogap}
\gapof{H}\;=\;\sup\set{\gamma\ge 0\,:\,\operatorname{spec}(H_\omega)\cap(0,\gamma)=\varnothing},
\end{equation}
the width of the spectral gap of $H_\omega$ above $0$ (with $\gapof{H}=0$ if no such
$\gamma>0$ exists). This is the quantity appearing in \eqref{eq:gap-def-intro}. When a
system has a uniform finite-volume gap---$\gapfin{\Lambda,m}(H)\ge\gamma$ for all
$\Lambda$ with $\gamma$ independent of $\Lambda$---the thermodynamic gap satisfies
$\gapof{H}\ge\gamma$; but the converse can fail, and a thermodynamic gap need not be
accompanied by a uniform finite-volume lower bound taken as an axiom. The
Nachtergaele--Sims--Young framework \cite{nsy-bulkgap} is designed precisely to work
with the infinite-volume gap directly rather than to assume uniform finite-size
bounds, and we adopt its stance.

\begin{remark}
\label{rem:cpw-finite}
Undecidability \cite{cpw-undecidable} is a statement about \eqref{eq:thermogap}, not
\eqref{eq:ngap}. For any fixed finite $\Lambda$ the number $\gapfin{\Lambda,n}(H)$ is
an eigenvalue difference of an explicit finite matrix and is computable to any
precision. What no algorithm can do is decide, from the local interaction data,
whether $\lim$-behavior produces $\gapof{H}>0$ or $\gapof{H}=0$. Every finite-volume
statement in this paper is therefore safe from the wall; every thermodynamic
statement is made \emph{conditional on} a standing gap hypothesis.
\end{remark}

\subsection{The uniformly gapped substack}
\label{sec:framework-substack}

\begin{definition}[Uniformly gapped substack]
\label{def:gap}
For a profinite probe $S$, an interaction family $H\in\Ham_{d,G}(S)$ is
\emph{uniformly gapped} if there exists $\gapbound>0$ with
$\einf_{s\in S}\gapof{H_s}\ge\gapbound$. The uniformly gapped systems form the
subpresheaf $\Gap_{d,G}\subseteq\Ham_{d,G}$ with
\[
\Gap_{d,G}(S)\;=\;\bigcup_{\gapbound>0}\Gap_{d,G}^{\ge\gapbound}(S),
\qquad
\Gap_{d,G}^{\ge\gapbound}(S)=\set{H\in\Ham_{d,G}(S)\,:\,\einf_{s\in S}\gapof{H_s}\ge\gapbound}.
\]
\end{definition}

The filtration by $\gapbound$ is genuine: the strata $\Gap_{d,G}^{\ge\gapbound}$ are
nested, $\Gap_{d,G}^{\ge\gapbound'}\subseteq\Gap_{d,G}^{\ge\gapbound}$ for
$\gapbound'\ge\gapbound$, and a family can belong to the union without belonging to
any single stratum uniformly as the probe varies. The following elementary
observation records the sheaf-theoretic shape of the object and will be sharpened,
conjecturally, in \Cref{sec:conjectures}.

\begin{proposition}[Restriction stability and the union structure]
\label{prop:union}
Let $S'\to S$ be a map of profinite sets and $H\in\Gap_{d,G}^{\ge\gapbound}(S)$. Then
the restriction $H|_{S'}$ lies in $\Gap_{d,G}^{\ge\gapbound}(S')$. Consequently each
$\Gap_{d,G}^{\ge\gapbound}$ is a subpresheaf of $\Ham_{d,G}$, and $\Gap_{d,G}$ is the
filtered union $\bigcup_{\gapbound>0}\Gap_{d,G}^{\ge\gapbound}$ of subpresheaves.
\end{proposition}

\begin{proof}
A map $g:S'\to S$ of profinite sets sends the family $s\mapsto H_s$ to
$s'\mapsto H_{g(s')}$. Since $g(S')\subseteq S$,
$\einf_{s'\in S'}\gapof{H_{g(s')}}\ge\einf_{s\in S}\gapof{H_s}\ge\gapbound$, so
$H|_{S'}\in\Gap_{d,G}^{\ge\gapbound}(S')$. Presheaf functoriality is immediate from
this together with the functoriality of $\Ham_{d,G}$. The union statement is
\Cref{def:gap}.
\end{proof}

\begin{remark}
\label{rem:pointwise-vs-uniform}
\Cref{prop:union} uses only $\einf$ monotonicity under restriction; it does
\emph{not} give the reverse implication, that a family gapped on every finite
quotient of $S$ is uniformly gapped on $S$. That reverse implication is a descent
statement, and whether it holds is the content of Conjecture III-2. The gap between
``gapped pointwise / at every finite resolution'' and ``uniformly gapped over the
whole profinite probe'' is exactly the phenomenon \Cref{def:gap} is built to track.
\end{remark}

\subsection{What ``open substack'' should mean}
\label{sec:framework-open}

Because $\BF$ is a Banach space, its condensation $\cond{\BF}$ is a condensed
$\R$-vector space and $\Ham_{d,G}$ sits over it. Openness of $\Gap_{d,G}$ inside
$\Ham_{d,G}$ can be phrased in the condensed setting as follows. Call a subobject
$U\subseteq\cond{\BF}$ \emph{condensed-open along a direction class $\mathcal{D}$} if
for every point $\Phi\in U$ there is $\varepsilon>0$ such that every $\Psi$ with
$\Psi-\Phi\in\mathcal{D}$ and $\norm{\Psi-\Phi}_{\Ffun}<\varepsilon$ again lies in
$U$, and if this holds compatibly for probes: pulling back along any $g:S\to\cond{\BF}$
whose image lies in the $\varepsilon$-ball around $\Phi$ lands in $U(S)$. For a full
topology on $\Ham_{d,G}$ one wants $\mathcal{D}$ to be all quasi-local directions and
$\varepsilon$ uniform; the honest situation, developed in \Cref{sec:stability}, is
that we can prove condensed-openness of $\Gap_{d,G}$ only \emph{on a stratum} and
only \emph{along the quasi-local perturbation class} for which the stability theorems
apply. Recording this scope precisely is not a weakness of the result but the exact
shape undecidability forces the result to take.

\section{The undecidability wall and the definitional stance}
\label{sec:wall}

\subsection{The Cubitt--P\'erez-Garc\'ia--Wolf theorem}
\label{sec:wall-cpw}

We restate the constraint precisely, in the form we will use.

\begin{theorem}[Cubitt--P\'erez-Garc\'ia--Wolf \cite{cpw-undecidable}]
\label{thm:cpw}
There is a fixed local Hilbert space dimension and an explicit, computable family
$n\mapsto H(n)$ of translation-invariant nearest-neighbour Hamiltonians on
$\Z^2$, indexed by a positive integer parameter $n$, with the following property.
The map that sends $n$ to the truth value of ``the thermodynamic-limit system
$H(n)$ is gapped'' is undecidable: no Turing machine, given $n$, halts with the
correct yes/no answer for all $n$. The problem is $\Pi_1$-hard, and the two cases are
sharply separated (a gap bounded below by a constant in the gapped case, dense
spectrum above the ground state in the gapless case), so the undecidability is not an
artifact of borderline or ambiguously-gapped systems.
\end{theorem}

Two features of \Cref{thm:cpw} matter for us. First, the hardness is already present
for translation-invariant, nearest-neighbour, two-dimensional systems---the most
classical and most physical corner of $\Ham_{d,G}$---so one cannot escape it by
restricting to a tame-looking subclass defined by locality or symmetry alone. Second,
the separation between the gapped and gapless cases is sharp, which forecloses the
hope that some ``soft'' or approximate criterion could decide the gap up to an
$\varepsilon$: the difficulty is not analytic delicacy near a threshold but genuine
computational undecidability.

The undecidability persists in one spatial dimension \cite{bcl-pg-1d-undecidable} and
under various symmetry restrictions, so it is not an artifact of high dimension. It
is, however, a \emph{fragile} phenomenon, in a way consonant with the viewpoint of
this paper: the undecidable families are exactly the ones that are not stably gapped.
Castilla-Castellano and Lucia \cite{castilla-lucia-instability} show that an
arbitrarily small local perturbation of the one-dimensional construction restores
decidability, by breaking the fine-tuned energy cancellation on which the halting
encoding depends. Undecidability lives on the boundary between gapped and gapless---%
precisely the locus the stability theory of \Cref{sec:stability} stays away from, and
precisely why that theory must assume, rather than certify, membership in the gapped
substack.

\subsection{Consequences for the substack}
\label{sec:wall-consequences}

\Cref{thm:cpw} has immediate and non-negotiable consequences for how $\Gap_{d,G}$ may
be studied.

\begin{corollary}[No uniform gap criterion]
\label{cor:no-criterion}
There is no algorithm that, given a finite description of a member
$H\in\Ham_{d,G}$ (its local interaction terms and symmetry data), decides whether
$H\in\Gap_{d,G}$ over the one-point probe. Equivalently, there is no computable
function $c:\Ham_{d,G}\to\R_{>0}$ such that $\gapof{H}\ge c(H)$ whenever $H$ is
gapped and $c(H)$ certifies gaplessness otherwise.
\end{corollary}

\begin{proof}
Such an algorithm, restricted to the CPW family $n\mapsto H(n)$ of \Cref{thm:cpw},
would decide gappedness of $H(n)$ from the finite data determining $H(n)$,
contradicting undecidability.
\end{proof}

Consequently, ``$H\in\Gap_{d,G}$'' cannot be the output of a general computation; it
is an input. We adopt the following stance, which we regard as forced rather than
chosen.

\begin{quote}
\textbf{Definitional stance.} The uniformly gapped substack $\Gap_{d,G}$ is
\emph{defined} by the gap hypothesis, not detected by a criterion. Membership is a
hypothesis one places on a family. The theorems of this paper take the form
``\emph{if} a family (or a stratum point) is uniformly gapped and satisfies stated
structural hypotheses, \emph{then} its neighborhood, correlations, and descent behavior
are controlled.'' No theorem asserts, for a general family, that the hypothesis holds.
\end{quote}

\begin{remark}
\label{rem:what-is-decidable}
It is worth marking the boundary precisely, because it is easy to overstate the wall
in either direction. \emph{Undecidable:} the thermodynamic gap of a general
translation-invariant family (\Cref{thm:cpw}). \emph{Decidable / computable:} the
finite-volume $n$-gap of any explicit finite $H_\Lambda$ (\Cref{rem:cpw-finite}); the
gap of any system for which one has a proof (frustration-free systems with a
Knabe-type or martingale bound, exactly solvable models, etc.); and (this is the
content of \Cref{sec:stability}) the \emph{persistence} of a gap under small
perturbations once a gap is known to exist with the right structural hypotheses.
Stability is a conditional, hypothesis-carrying statement, and conditional statements
are exactly what undecidability permits.
\end{remark}

\subsection{Why the program survives the wall}
\label{sec:wall-survives}

It might seem that undecidability is fatal to a program that puts $\Gap_{d,G}$ at its
center. It is not, for the same reason that undecidability of the halting problem is
not fatal to the theory of computation: one develops the theory of the objects
\emph{assuming} the property, and proves structural and stability results about the
class so defined. The moduli-theoretic viewpoint is well suited to this. A stack does
not need a decision procedure for membership to be a useful object; it needs good
functorial behavior, a sensible topology, and stability of its distinguished
subobjects under the maps one cares about. Those are precisely what Theorems
III-A--III-C provide, on the stratum where they can be proved. The undecidability
wall bounds the \emph{global} reach of the theory (one cannot hope for a single
computable classification of all gapped systems), but it leaves the local and
stratified theory intact, and it is the local and stratified theory that feeds Parts
IV and V\@.

\section{Finite-volume gap continuity}
\label{sec:finite}

We begin with the one unconditional result: at finite volume, the gap is a Lipschitz
function of the interaction, and the gapped locus is open. This is the finite-volume
shadow of the openness we would like for $\Gap_{d,G}$, and it is the part of the
picture undecidability leaves entirely alone.

\subsection{Weyl continuity of eigenvalues}
\label{sec:finite-weyl}

\begin{lemma}[Weyl monotonicity]
\label{lem:weyl}
Let $A,B$ be self-adjoint operators on a finite-dimensional Hilbert space with
eigenvalues $\lambda_0(A)\le\cdots\le\lambda_{N-1}(A)$ and
$\lambda_0(B)\le\cdots\le\lambda_{N-1}(B)$ in non-decreasing order. Then for every
$k$,
\[
\abs{\lambda_k(A)-\lambda_k(B)}\;\le\;\norm{A-B},
\]
where $\norm{\cdot}$ is the operator norm.
\end{lemma}

\begin{proof}
This is Weyl's inequality. By the min--max theorem,
\[
\lambda_k(A)=\min_{\dim V=k+1}\ \max_{0\ne v\in V}\frac{\ip{v,Av}}{\ip{v,v}}.
\]
For any subspace $V$ and unit vector $v\in V$,
$\ip{v,Av}=\ip{v,Bv}+\ip{v,(A-B)v}\le\ip{v,Bv}+\norm{A-B}$, so
$\lambda_k(A)\le\lambda_k(B)+\norm{A-B}$. Symmetry in $A,B$ gives the reverse, and the
two together give the claim.
\end{proof}

To pass from operator norm to the interaction $\Ffun$-norm we use finite-volume norm
equivalence, which is the elementary end of the Lieb--Robinson machinery of Part I\@.

\begin{lemma}[Finite-volume norm domination]
\label{lem:normdom}
Fix a finite $\Lambda\subset L$. There is a constant $C_\Lambda<\infty$, depending on
$\Lambda$ and the $\Ffun$-function, such that for all interactions $\Phi,\Psi$,
\[
\norm{H_\Lambda(\Phi)-H_\Lambda(\Psi)}\;\le\;C_\Lambda\,\norm{\Phi-\Psi}_{\Ffun}.
\]
\end{lemma}

\begin{proof}
$H_\Lambda(\Phi)-H_\Lambda(\Psi)=\sum_{X\subseteq\Lambda}\big(\Phi(X)-\Psi(X)\big)$,
a finite sum of at most $2^{\abs{\Lambda}}$ terms. The $\Ffun$-norm dominates the
per-term operator norm up to the reweighting factor built into the norm
\cite{nsy-quasilocality-1}; summing the finitely many reweighting factors over
$X\subseteq\Lambda$ yields a finite $C_\Lambda$ with the stated bound.
\end{proof}

\subsection{Theorem III-A}
\label{sec:finite-thmA}

\begin{thmA}[Finite-volume gap continuity and finite-resolution openness]
\label{thm:IIIA}
Fix a finite $\Lambda\subset L$ and an integer $n\ge 1$.
\begin{enumerate}
\item[(i)] The $n$-gap is $2C_\Lambda$-Lipschitz in the interaction norm:
\[
\big\lvert\gapfin{\Lambda,n}(\Phi)-\gapfin{\Lambda,n}(\Psi)\big\rvert
\;\le\;2\,C_\Lambda\,\norm{\Phi-\Psi}_{\Ffun}
\qquad\text{for all }\Phi,\Psi.
\]
In particular $\Phi\mapsto\gapfin{\Lambda,n}(\Phi)$ is continuous on $\BF$.
\item[(ii)] For every $\delta>0$ the finite-volume gapped locus
$\set{\Phi\in\BF\,:\,\gapfin{\Lambda,n}(\Phi)>\delta}$ is open in $\BF$.
\item[(iii)] Along any $S$-continuous family $H\in\Ham_{d,G}(S)$ over a profinite
probe $S$, the map $s\mapsto\gapfin{\Lambda,n}(H_s)$ is continuous, and for every
$\delta>0$ the set $\set{s\in S\,:\,\gapfin{\Lambda,n}(H_s)>\delta}$ is open in $S$.
\end{enumerate}
\end{thmA}

\begin{proof}
(i) By \eqref{eq:ngap},
$\gapfin{\Lambda,n}(\Phi)-\gapfin{\Lambda,n}(\Psi)
=\big(\lambda_n(H_\Lambda(\Phi))-\lambda_n(H_\Lambda(\Psi))\big)
-\big(\lambda_0(H_\Lambda(\Phi))-\lambda_0(H_\Lambda(\Psi))\big)$.
\Cref{lem:weyl} bounds each parenthesized difference by
$\norm{H_\Lambda(\Phi)-H_\Lambda(\Psi)}$, and \Cref{lem:normdom} bounds that by
$C_\Lambda\norm{\Phi-\Psi}_{\Ffun}$. The triangle inequality gives the factor $2$.

(ii) Immediate from (i): a $2C_\Lambda$-Lipschitz function has open strict
super-level sets.

(iii) An $S$-continuous family is a continuous map $S\to\BF$; composing with the
continuous $\gapfin{\Lambda,n}$ of (i) gives a continuous $S\to\R$, and continuity
gives openness of super-level sets.
\end{proof}

\begin{remark}[Why the fixed integer $n$]
\label{rem:why-n}
If one used the ``physical'' gap $\lambda_m-\lambda_0$ with $m$ the ground-state
multiplicity (itself a function of $\Phi$), the map $\Phi\mapsto$ gap could jump
where $m$ changes, and Lipschitz continuity would fail there. Fixing $n$ removes the
discontinuity by decoupling the level index from the multiplicity. On any region
where the ground-state multiplicity is locally constant and equal to $m$, the
physical gap agrees with the $m$-gap and Theorem III-A applies verbatim. The
transverse-field Ising computation of \Cref{sec:numerics} exhibits exactly this: in
the ferromagnetic phase the two lowest levels form a near-degenerate ground sector
($m=2$ in the thermodynamic limit), and the physically meaningful gap is the $2$-gap,
not the $1$-gap. As $\Lambda$ grows, the ground-state multiplicity can change at
isolated interaction values (level crossings or accidental degeneracies); the
fixed-$n$ device keeps $\gapfin{\Lambda,n}$ a well-defined, continuous
ordered-eigenvalue difference across such events, and one reads off the physical gap
by taking $n$ equal to the ground-sector dimension on the region of interest, where it
is locally constant. The continuity of \emph{each} ordered eigenvalue (\Cref{lem:weyl})
never fails; only the identification of which level index is ``the gap'' can shift, and
that shift happens on a measure-zero set of couplings.
\end{remark}

\subsection{Semicontinuity of the thermodynamic gap}
\label{sec:finite-semicont}

At the thermodynamic level the gap is only \emph{lower} semicontinuous in general,
and even that requires care about the mode of convergence. We record the statement we
can make and mark the pitfall.

\begin{proposition}[Conditional persistence and one-sided continuity of the gap]
\label{prop:lsc}
Suppose $\Phi_j\to\Phi$ in $\BF$ and that all $\Phi_j$ and $\Phi$ admit
infinite-volume ground states obtained as weak-$*$ limits of finite-volume ground
states in the sense of Part II\@. If, for a subsequence, the GNS gaps satisfy
$\gapof{\Phi_j}\ge\gamma$ for a fixed $\gamma>0$, then $\gapof{\Phi}\ge\gamma$
provided the ground states converge weak-$*$ and the gap is realized by a convergent
family of low-lying states. This $\gamma$-persistence is the conditional
\emph{upper}-semicontinuity direction. Unconditionally the thermodynamic gap is only
\emph{lower} semicontinuous, $\gapof{\Phi}\le\liminf_j\gapof{\Phi_j}$, which supplies
no lower bound (the collapse $\gapof{\Phi_j}=c>0\to\gapof{\Phi}=0$ satisfies it);
without the auxiliary hypotheses no lower bound on $\gapof{\Phi}$ survives at all.
\end{proposition}

\begin{proof}[Proof sketch]
The argument is the standard one for persistence of a spectral gap under strong
resolvent convergence, adapted to the GNS setting of \cite{nsy-bulkgap}: a spectral
gap is the statement that the ground state minimizes energy with a margin
$\gamma$ against all orthogonal excitations, an inequality
$\ip{\psi,(H_\omega-\gamma)\psi}\ge 0$ for $\psi\perp\Omega_\omega$; such inequalities
pass to weak-$*$ limits of states when the energy functionals converge, which they do
along $\BF$-convergent interactions by the Lieb--Robinson continuity of the dynamics
(Part I)\@. Failure of the auxiliary hypotheses---non-convergence of ground states, or
a low-lying spectrum that ``escapes'' in the limit---breaks the inequality; we do not
claim a bound in that case, and by \Cref{cor:no-criterion} no unconditional bound can
exist.
\end{proof}

\begin{remark}
\label{rem:lsc-caveat}
\Cref{prop:lsc} is deliberately hedged. It would be an error, and a violation of the
undecidability wall, to state ``the thermodynamic gap is continuous in the
interaction'' as an unconditional theorem: continuity would let one certify gaps by
approximation, which \Cref{cor:no-criterion} forbids. What fails is \emph{upper}
semicontinuity: along a sequence of gapped interactions a low-lying excited level can
descend to the ground-state energy in the limit, collapsing the gap discontinuously,
and this is precisely the mechanism the undecidable families of \cite{cpw-undecidable}
exploit: the gap is present at every finite approximation and disappears only in the
thermodynamic limit. Lower semicontinuity under convergence of the low-lying spectral
data is therefore the honest statement, and it is exactly what the stability theory of
\Cref{sec:stability} upgrades, on its stratum, to genuine openness with a uniform
radius.
\end{remark}

\section{Stability on the frustration-free / LTQO stratum}
\label{sec:stability}

We now reach the heart of the paper: the recasting of the
Bravyi--Hastings--Michalakis, Michalakis--Zwolak and Nachtergaele--Sims--Young
stability theorems as a statement that $\Gap_{d,G}$ contains condensed-open
neighborhoods on a well-defined stratum. Throughout, we are scrupulous about
hypotheses; the stratum is exactly where the hypotheses hold with constants uniform
in the system size.

\subsection{Frustration-freeness and local topological quantum order}
\label{sec:stability-hypotheses}

\begin{definition}[Frustration-free interaction]
\label{def:ff}
An interaction $\Phi$ is \emph{frustration-free} if there is a decomposition
$H_\Lambda(\Phi)=\sum_{x\in\Lambda}h_x$ with each $h_x\ge 0$ a local term, such that
the finite-volume ground-state space $\gsp_\Lambda=\ker H_\Lambda(\Phi)$ is the common
kernel $\bigcap_x\ker h_x$; equivalently the finite-volume ground-state energy is $0$
and every ground state is annihilated by every local term.
\end{definition}

Frustration-freeness is a genuine restriction. Generic gapped systems are not
frustration-free; the class includes the paradigmatic exactly-solvable models
(Kitaev's toric code and honeycomb model \cite{kitaev-honeycomb}, AKLT-type chains,
group-cohomology fixed-point models \cite{cglw-cohomology}) and the parent
Hamiltonians of matrix-product and PEPS states, but it excludes, for instance,
generic Chern insulators. We never extend a frustration-free theorem to a
non-frustration-free system.

\begin{definition}[Local topological quantum order, LTQO]
\label{def:ltqo}
A frustration-free family with ground-state projections $P_\Lambda$ onto
$\gsp_\Lambda$ satisfies \emph{LTQO} with decay $\Omega_{\mathrm{LTQO}}(\cdot)$ if for
every ball $B_r(x)$ of radius $r$ and every observable $O$ supported on $B_r(x)$,
\[
\norm{P_\Lambda\,O\,P_\Lambda-\frac{\Tr(P_{B_{r'}(x)}O)}{\Tr P_{B_{r'}(x)}}\,P_\Lambda}
\;\le\;\norm{O}\,\Omega_{\mathrm{LTQO}}(r'-r)
\]
for $r<r'$, where $\Omega_{\mathrm{LTQO}}$ is a fixed rapidly-decaying function
independent of $\Lambda$. Informally: ground states are locally indistinguishable, up
to a boundary correction that decays in the distance to the region's edge.
\end{definition}

The LTQO condition is the operator-algebraic expression of ``topological'' ground-state
degeneracy: the degenerate ground states cannot be told apart by any local
measurement. Michalakis and Zwolak \cite{michalakis-zwolak} showed that, together with
a uniform local gap, LTQO is essentially equivalent to stability of the gap and to an
area law for the ground-state entanglement.

\begin{definition}[The frustration-free / LTQO stratum]
\label{def:stratum}
Fix $\gamma>0$, a decay function $\Omega_{\mathrm{LTQO}}$, an $\Ffun$-function, and
constants controlling the local-term norms and interaction range. The
\emph{frustration-free / LTQO stratum}
$\mathfrak{S}=\mathfrak{S}(\gamma,\Omega_{\mathrm{LTQO}},\Ffun)\subseteq\Ham_{d,G}$
is the collection of interactions $\Phi$ that are frustration-free
(\Cref{def:ff}), have finite-volume gap $\gapfin{\Lambda,m_\Lambda}(\Phi)\ge\gamma$
uniformly in $\Lambda$ (with $m_\Lambda=\dim\gsp_\Lambda$), and satisfy LTQO with
decay $\Omega_{\mathrm{LTQO}}$, all with the fixed constants. As a subpresheaf,
$\mathfrak{S}(S)$ consists of $S$-families landing in $\mathfrak{S}$ for every $s\in S$
with the constants uniform over $S$.
\end{definition}

By construction $\mathfrak{S}\subseteq\Gap_{d,G}$: the uniform finite-volume gap
$\gamma$ forces $\gapof{\Phi}\ge\gamma$ (\Cref{sec:framework-gap}). The stratum is the
domain of the stability theorems.

\begin{remark}[The stratum is a proper subclass; what lies outside it]
\label{rem:stratum-boundary}
Frustration-freeness (\Cref{def:ff}) is essential to the arguments of
\Cref{thm:bhm,thm:mz,thm:nsy}: the relative-bound and martingale/telescoping
estimates that produce the size-independent threshold rest on the common-kernel
structure and the local projectors $P_\Lambda$, and on LTQO stated through those
projectors. Many physically important gapped systems are \emph{not} of this form.
Free-fermion topological insulators and Chern insulators, for instance, are gapped
and stable, but they are not frustration-free commuting-projector systems (indeed
the chiral ones cannot be, by the Kapustin--Fidkowski obstruction \cite{kapustin-fidkowski} discussed in
\Cref{sec:disc-walls}), so they lie outside $\mathfrak{S}$, and Theorem III-B says
nothing about them. Their stability is real but is established by different
technology (single-particle / spectral-localizer and $K$-theoretic methods,
\cite{prodan-sb}), not by the frustration-free stability theorems recast here. We do
not attempt to widen $\mathfrak{S}$ to cover them; whether a single condensed-openness
statement subsumes both classes is part of the open Conjecture III-1, not something
the present theorem delivers.
\end{remark}

\subsection{The stability theorems, as cited}
\label{sec:stability-cited}

We quote the results we use in the form relevant here; the constants are those of the
original sources.

\begin{theorem}[Bravyi--Hastings--Michalakis \cite{bhm-stability}; short proof
\cite{bravyi-hastings-shortproof}]
\label{thm:bhm}
Let $\Phi$ be a frustration-free interaction satisfying the topological-order
conditions TQO-1 and TQO-2 (LTQO) with a uniform local gap $\gamma>0$. There exists
$\varepsilon_0>0$, depending only on $\gamma$, the LTQO decay, the local Hilbert space
dimension, and the lattice geometry---and \emph{not} on the system size---such that
for any perturbation $V=\sum_x V_x$ with $\sup_x\norm{V_x}\le\varepsilon<\varepsilon_0$
and sufficiently fast decay, the perturbed finite-volume Hamiltonians
$H_\Lambda(\Phi)+V_\Lambda$ have a spectral gap above their ground-state sector
bounded below by $\gamma/2$ (say), uniformly in $\Lambda$, with the ground-state
splitting within the sector decaying faster than any polynomial in the system size.
\end{theorem}

\begin{theorem}[Michalakis--Zwolak \cite{michalakis-zwolak}]
\label{thm:mz}
For frustration-free Hamiltonians, LTQO together with a uniform local gap implies
stability of the gap under quasi-local perturbations, with a size-independent
threshold; conversely, stability plus an area law essentially forces LTQO\@. The
perturbed system retains a uniform gap and a ground-state space continuously connected
to the unperturbed one by a spectral flow.
\end{theorem}

\begin{theorem}[Nachtergaele--Sims--Young \cite{nsy-bulkgap}]
\label{thm:nsy}
For frustration-free, topologically ordered quantum lattice systems, the bulk gap in
the infinite-volume GNS representation is stable under quasi-local perturbations
controlled by an $\Ffun$-function, with a threshold independent of the finite-volume
approximations. The result is proved directly for the infinite-volume gap
\eqref{eq:thermogap}, without assuming a uniform finite-size lower bound as a separate
axiom.
\end{theorem}

The three results are complementary: \cite{bhm-stability,bravyi-hastings-shortproof}
establish the finite-volume, size-independent threshold; \cite{michalakis-zwolak}
identifies LTQO as the precise condition and connects it to the area law;
\cite{nsy-bulkgap} lifts the conclusion to the thermodynamic gap directly. What they
share, and what we need, is a \emph{single} quantity: a perturbation radius
$\varepsilon_0>0$ that does not shrink with system size.

\subsection{Theorem III-B: stability as condensed-openness}
\label{sec:stability-thmB}

\begin{thmB}[Conditional condensed-openness of $\Gap_{d,G}$ on the stratum]
\label{thm:IIIB}
Let $\mathfrak{S}=\mathfrak{S}(\gamma,\Omega_{\mathrm{LTQO}},\Ffun)$ be a
frustration-free / LTQO stratum (\Cref{def:stratum}), and let $\varepsilon_0>0$ be the
size-independent perturbation threshold of \Cref{thm:bhm,thm:mz,thm:nsy}, depending
only on $\gamma$, the LTQO decay $\Omega_{\mathrm{LTQO}}$, the $\Ffun$-function, and
$d$. Then for every
$\Phi\in\mathfrak{S}$ and every quasi-local perturbation direction
$V=\sum_x V_x$ with $\sup_x\norm{V_x}\le 1$ and $\Ffun$-controlled decay, the segment
\[
\set{\Phi+tV\,:\,0\le t<\varepsilon_0}\ \subseteq\ \Gap_{d,G},
\]
with thermodynamic gap bounded below by $\gamma/2$ throughout.\par\noindent
Consequently $\Gap_{d,G}$ is condensed-open along the quasi-local direction class at
each point of
$\mathfrak{S}$: for every profinite probe $S$ and every $S$-family $H$ landing in the
$\varepsilon_0$-ball of quasi-local perturbations around a point of
$\mathfrak{S}$, one has $H\in\Gap_{d,G}^{\ge\gamma/2}(S)$. In words: on the
frustration-free / LTQO stratum, the uniformly gapped substack contains a
condensed-open neighborhood of each point along quasi-local perturbation directions.
\end{thmB}

\begin{proof}
Fix $\Phi\in\mathfrak{S}$ and a quasi-local direction $V$ normalized as stated. For
$0\le t<\varepsilon_0$, the perturbation $tV$ has per-site strength
$\sup_x\norm{tV_x}\le t<\varepsilon_0$ and the $\Ffun$-controlled decay required by
\Cref{thm:bhm,thm:mz,thm:nsy}. Since $\Phi$ lies in the stratum, it satisfies the
frustration-freeness, uniform local gap $\gamma$, and LTQO hypotheses of those
theorems with the fixed constants; therefore each theorem applies to $\Phi+tV$ and
yields a gap above the ground-state sector bounded below by $\gamma/2$, uniformly in
the finite volume, and---by \Cref{thm:nsy}---a thermodynamic gap
$\gapof{\Phi+tV}\ge\gamma/2$. Since $\gamma/2>0$ is independent of $t$ in the range,
the whole segment lies in $\Gap_{d,G}^{\ge\gamma/2}\subseteq\Gap_{d,G}$.

For the condensed statement, let $S$ be a profinite probe and $H\in\Ham_{d,G}(S)$ a
family with $H_s=\Phi+t(s)V(s)$ where, for each $s$, $\Phi\in\mathfrak{S}$, $V(s)$ is a
normalized quasi-local direction, and $t(s)<\varepsilon_0$, with all constants uniform
over $S$ (this is what it means for the family to land in the $\varepsilon_0$-ball
along quasi-local directions with stratum-uniform data). The pointwise bound just
proved gives $\gapof{H_s}\ge\gamma/2$ for every $s\in S$, hence
$\einf_{s\in S}\gapof{H_s}\ge\gamma/2$ and $H\in\Gap_{d,G}^{\ge\gamma/2}(S)$ by
\Cref{def:gap}. This is condensed-openness along the quasi-local class at $\Phi$ in
the sense of \Cref{sec:framework-open}.
\end{proof}

\begin{remark}[Exactly what is and is not claimed]
\label{rem:IIIB-scope}
Theorem III-B is a stability statement and nothing more. It does \emph{not} decide
whether any given $\Phi$ belongs to $\mathfrak{S}$: that requires knowing $\Phi$ is
frustration-free with a uniform gap $\gamma$, which is an input, consistent with
\Cref{cor:no-criterion}. It does \emph{not} assert openness of $\Gap_{d,G}$ at points
outside $\mathfrak{S}$, nor along non-quasi-local directions, nor with a threshold
uniform over all of $\Gap_{d,G}$. The three ``not''s are not timidity; each is a place
where a stronger claim would either be false or would contradict undecidability. The
uniform-in-size threshold $\varepsilon_0$ is the entire mechanism: it is what turns a
family of finite-volume openness statements (Theorem III-A, which has $C_\Lambda\to\infty$
as $\Lambda$ grows) into a single thermodynamic openness statement with a radius that
does not collapse.
\end{remark}

\begin{corollary}[Openness of the stratum stratum-wise]
\label{cor:stratum-open}
The intersection $\Gap_{d,G}\cap(\Phi+\varepsilon_0\,\mathcal{D}_{\mathrm{ql}})$
contains the full $\varepsilon_0$-ball along the quasi-local direction class
$\mathcal{D}_{\mathrm{ql}}$ at each $\Phi\in\mathfrak{S}$. In particular the assignment
$\Phi\mapsto\gamma/2$ is a lower bound for the gap that is \emph{locally constant along
quasi-local perturbations} on $\mathfrak{S}$, which is the local constancy the phase
label $\nu_f$ of the program requires.
\end{corollary}

\begin{proof}
Immediate from Theorem III-B: every point of the ball has gap $\ge\gamma/2$, so the
lower bound $\gamma/2$ is constant on the ball, and local constancy of a lower bound
is what a locally constant phase label needs to be well defined on the gapped locus.
\end{proof}

\subsection{Spectral flow and automorphic equivalence}
\label{sec:stability-flow}

The stability of \Cref{sec:stability-thmB} comes with a dynamical companion: within a
uniformly gapped neighborhood the ground-state spaces are related by a quasi-local
automorphism, the spectral flow. This is the rigorous meaning of ``the same phase''
along a gapped path, and it is the interface with Part I and with the equivalence
class $\Weq$.

\begin{proposition}[Spectral flow inside a gapped segment]
\label{prop:flow}
Let $t\mapsto\Phi+tV$, $0\le t\le t_1<\varepsilon_0$, be a segment as in Theorem III-B,
so $\gapof{\Phi+tV}\ge\gamma/2$ throughout. Then there is a quasi-local, $\Ffun$-norm
continuous cocycle of $*$-automorphisms $(\alpha_{t})_{0\le t\le t_1}$ of the
quasi-local algebra, generated by a time-dependent quasi-local interaction, such that
$\alpha_t$ maps the ground-state space of $\Phi$ to that of $\Phi+tV$, and
$\alpha_t\to\mathrm{id}$ as $\gamma\to\infty$ in the sense of the Lieb--Robinson
tails. The automorphisms $\alpha_t$ belong to the class $\Weq$ of gapped quasi-local
equivalences.
\end{proposition}

\begin{proof}[Proof sketch]
This is the quasi-adiabatic continuation / spectral-flow construction of Hastings and
Wen \cite{hastings-wen-qac} in the quasi-local formulation of Bachmann, Michalakis,
Nachtergaele and Sims \cite{bmns-automorphic} and Nachtergaele, Sims and Young
\cite{nsy-quasilocality-2}. Given a uniformly gapped path $\Phi+tV$, one forms the
spectral-flow generator $D_t=\int F_\gamma(s)\,\tau^{(t)}_s\!\big(\tfrac{d}{dt}H(t)\big)\,ds$
with a filter $F_\gamma$ whose Fourier support avoids the gap; the gap lower bound
$\gamma/2$ makes $D_t$ quasi-local with $\Ffun$-tails controlled by $\gamma$, and the
resulting flow $\alpha_t$ intertwines the ground-state spaces. Continuity in the
$\Ffun$-norm and membership in $\Weq$ are then as in Part I\@. The $\gamma\to\infty$
limit sharpens the filter to a delta and the flow to the identity on the ground-state
sector.
\end{proof}

\begin{remark}
\label{rem:flow-W}
\Cref{prop:flow} is what makes ``$\Gap_{d,G}$ modulo $\Weq$'' a sensible quotient
along stratum segments: adiabatically connected uniformly gapped systems are
$\Weq$-equivalent by an explicit quasi-local automorphism. Conjecture III-3 below
proposes that this is the whole story (that the internal path components of
$\Gap_{d,G}$ under quasi-adiabatic continuation are exactly the operator-algebraic
gapped phases), but the general statement is beyond what \Cref{prop:flow} proves,
because a general uniformly gapped path need not lie in a single stratum.
\end{remark}

\section{Uniform clustering from a uniform gap}
\label{sec:clustering}

The third theorem is structural: it isolates a property that \emph{uniformly} gapped
families have and pointwise-gapped families need not. A gap forces exponential decay
of correlations; a \emph{uniform} gap forces the decay rate to be uniform over the
base. This is the sense in which $\Gap_{d,G}$, and not the naive pointwise-gapped
locus, is the correct object.

\subsection{Gap implies clustering}
\label{sec:clustering-gap}

\begin{theorem}[Hastings--Koma \cite{hastings-koma}; Nachtergaele--Sims
\cite{nachtergaele-sims-clustering}]
\label{thm:hk}
Let $H$ have a thermodynamic gap $\gapof{H}\ge\gamma>0$ above a ground state
$\omega$, and let the dynamics obey a Lieb--Robinson bound with velocity $v$ and
$\Ffun$-tails. Then there are constants $c,C>0$, depending only on $\gamma$, $v$ and
the $\Ffun$-function, such that for local observables $A,B$ with disjoint supports,
\[
\big\lvert\omega(AB)-\omega(A)\,\omega(B)\big\rvert
\;\le\;C\,\norm{A}\,\norm{B}\;e^{-c\,\dist(\supp A,\supp B)},
\]
with $c\sim\gamma/v$ up to the geometric factors. The correlation length is
$\xi=c^{-1}$.
\end{theorem}

The precise exponent and constants are as in \cite{hastings-koma}; the
Lieb--Robinson input is the exponential-clustering theorem of
\cite{nachtergaele-sims-clustering}. We use them as cited.

\subsection{Theorem III-C: uniformity over the base}
\label{sec:clustering-thmC}

\begin{thmC}[Uniform clustering on a uniformly gapped family]
\label{thm:IIIC}
Let $H\in\Gap_{d,G}^{\ge\gapbound}(S)$ be a uniformly gapped family over a profinite
probe $S$, so $\einf_{s\in S}\gapof{H_s}\ge\gapbound>0$, and suppose the family has a
uniform Lieb--Robinson velocity $v$ and common $\Ffun$-function (as provided by Part I
for an $S$-continuous family in $\BF$). Then there are constants $c,C>0$, depending
only on $\gapbound$, $v$ and $\Ffun$---and \emph{not} on $s$---such that for every
$s\in S$ and all local $A,B$ with disjoint supports,
\[
\big\lvert\omega_s(AB)-\omega_s(A)\,\omega_s(B)\big\rvert
\;\le\;C\,\norm{A}\,\norm{B}\;e^{-c\,\dist(\supp A,\supp B)}.
\]
In particular the correlation length function $s\mapsto\xi(s)$ is bounded above by
$c^{-1}$ uniformly over the profinite base.
\end{thmC}

\begin{proof}
Apply \Cref{thm:hk} at each $s\in S$. Its constants $c,C$ depend on the input data
$(\gamma,v,\Ffun)$ only. By hypothesis the gap input can be taken to be the uniform
value $\gamma=\gapbound$ for every $s$ (since $\gapof{H_s}\ge\gapbound$), and the
velocity and $\Ffun$-function are the common ones supplied by Part I for the
$S$-continuous family. Thus the constants produced by \Cref{thm:hk} are the same for
every $s$, and the clustering bound holds with $s$-independent $c,C$. The bound on
$\xi(s)=c^{-1}$ is the final clause.
\end{proof}

\begin{corollary}[Failure for pointwise-gapped families]
\label{cor:pointwise-fail}
If instead $H_s$ is gapped for every $s\in S$ but $\einf_s\gapof{H_s}=0$, then the
per-point clustering rate $c(s)\sim\gapof{H_s}/v$ degenerates to $0$ along a sequence
$s_j$ with $\gapof{H_{s_j}}\to 0$, and no uniform correlation-length bound exists. The
uniform-gap hypothesis of Theorem III-C is therefore necessary for the conclusion, not
merely sufficient.
\end{corollary}

\begin{proof}
The Hastings--Koma rate is proportional to the gap; along $s_j$ it tends to $0$, so
$\xi(s_j)\to\infty$ and $\sup_s\xi(s)=\infty$.
\end{proof}

\begin{remark}
\label{rem:clustering-meaning}
\Cref{cor:pointwise-fail} is the technical vindication of \Cref{def:gap}. It shows
that the difference between ``gapped at every probe point'' and ``uniformly gapped''
is not a formality: only the latter guarantees a uniform correlation length, and a
uniform correlation length is what one needs for the area laws, split property, and
locally computable invariants that Parts II and IV build on. The $\einf_s$ in the
definition of $\Gap_{d,G}$ is precisely the hypothesis that rules out
\Cref{cor:pointwise-fail}.
\end{remark}

\section{The structure of the gapped substack: conjectures}
\label{sec:conjectures}

Beyond the stratum, we do not have theorems; we have conjectures, and we state them
as such, with the reasons they are plausible and the obstacles to proving them. Each
is scoped to respect the undecidability wall.

\subsection{Openness in the full condensed topology}
\label{sec:conj-open}

\begin{conjone}[Openness on the physical stratum; finite-quotient detection]
\label{conj:III1}
On the physically relevant stratum---the closure, in a suitable sense, of the
frustration-free / LTQO systems under quasi-local gapped equivalence, together with
the systems for which a stability theory holds with uniform constants---$\Gap_{d,G}$
is an \emph{open} condensed substack of $\Ham_{d,G}$. Moreover the uniform-gap sheaf
condition of \Cref{def:gap} is detected by finite quotients of profinite probes: a
family over $S=\varprojlim S_i$ lies in $\Gap_{d,G}(S)$ if and only if its pushforwards
to the finite quotients $S_i$ lie in the corresponding finite-resolution gapped loci
with a common bound $\gapbound$.
\end{conjone}

\emph{Evidence and obstruction.} Theorem III-B is exactly this statement restricted to
$\mathfrak{S}$ and to quasi-local directions; the conjecture asserts that the phenomenon
persists on the larger physical stratum and in all directions. The obstruction to a
proof is precisely undecidability: openness in the \emph{full} $\Ham_{d,G}$, with a
uniform radius, would let one certify gaps of nearby systems, and near a CPW-type
embedding \cite{cpw-undecidable} the gap can turn on and off in a way no uniform radius
can survive. Restricting to the physical stratum is what excludes those pathological
directions; making ``physical stratum'' precise, so that the restriction is both
honest and non-vacuous, is the mathematical work the conjecture packages. The
finite-quotient detection clause is a descent statement and is closely tied to
Conjecture III-2.

\subsection{Uniform-gap descent along profinite covers}
\label{sec:conj-descent}

The disorder-hull thread of the program (\cite{bellissard-ncg-qhe,kellendonk-tilings}
and Part II) makes the profinite structure of $S$ physical: for a finite
local-configuration set $Q$ (the letter $Q$ denotes the disorder alphabet, kept
deliberately distinct from the locality $\Ffun$-function of Part I) the hull
$\dis=Q^{\Z^d}$ is compact and totally
disconnected, hence profinite, and a disordered family $\omega\mapsto H_\omega$ is
literally a $\dis$-point of $\Ham_{d,G}$. Finite quotients of $\dis$ are
finite-resolution disorder data.

\begin{conjtwo}[Uniform-gap descent]
\label{conj:III2}
Let $\dis=Q^{\Z^d}=\varprojlim_i\dis_i$ be a disorder hull with its profinite
presentation, and let $H\in\Ham_{d,G}(\dis)$ be a disordered family. Then $H$ is
uniformly gapped, $H\in\Gap_{d,G}(\dis)$, if and only if each finite-resolution
approximation $H^{(i)}\in\Ham_{d,G}(\dis_i)$ is gapped with a common bound $\gapbound$
independent of $i$. Equivalently, uniform gappedness is a \emph{closed} condition on
the inverse system $(\dis_i)_i$: the uniformly gapped families are the inverse limit
of the finite-resolution uniformly-gapped ones at a fixed $\gapbound$.
\end{conjtwo}

\emph{Evidence and obstruction.} The ``only if'' direction is \Cref{prop:union}: a
uniform gap over $\dis$ restricts to a uniform gap over each $\dis_i$. The content is
the ``if'' direction---that compatible finite-resolution gaps at a common $\gapbound$
assemble to a uniform gap over the limit. This is a genuine descent question. The
optimistic reason to expect it: continuity of the gap along $\BF$-convergent families
(Theorem III-A at each finite volume) plus compactness of $\dis$ suggests that a
common finite-resolution bound should survive the limit. The obstruction: the gap is a
\emph{thermodynamic}, not finite-volume, quantity, and \Cref{prop:lsc} shows only
lower semicontinuity in general; the descent could fail if the low-lying spectrum
reorganizes in the inverse limit. The conjecture asserts it does not, on the disorder
hulls, at fixed $\gapbound$. A proof would presumably route through the crossed-product
algebra $C(\dis)\rtimes\Z^d$ of Part II and a gap-labelling argument in the manner of
Bellissard \cite{bellissard-ncg-qhe}.

\subsection{Quasi-adiabatic continuation realizes the equivalence class}
\label{sec:conj-qac}

\begin{conjthree}[QAC path components are the operator-algebraic phases]
\label{conj:III3}
\leavevmode\par\noindent
Quasi-adiabatic continuation defines, on $\Gap_{d,G}$, the internal path components
realizing the equivalence class $\Weq$: two uniformly gapped families are
$\Weq$-equivalent if and only if they are connected by a uniformly gapped path along
which the spectral flow of \Cref{prop:flow} is quasi-local. Consequently
\[
\pi_0\Shape\!\big(\Gap_{d,G}[\Weq^{-1}]\big)\;\cong\;
\set{\text{operator-algebraic gapped ground-state phases}}
\]
of Ogata \cite{ogata-classification-review}, and the boxed slogan of the program---%
\emph{a topological phase is a component of the stabilized condensed stack of gapped
systems}---is, on this identification, a theorem rather than a definition.
\end{conjthree}

\emph{Evidence and obstruction.} \Cref{prop:flow} proves one direction locally: a
uniformly gapped segment in a stratum gives a quasi-local $\Weq$-equivalence. The
theory of automorphic equivalence in \cite{bmns-automorphic} and the classification
program of \cite{ogata-classification-review} make the target set precise in $d=1$
and, for on-site finite symmetry, in $d=2$. The obstruction to the full statement is
twofold: a general uniformly gapped path need not remain in a single stratum, so
\Cref{prop:flow} does not immediately apply along it; and the identification of
$\pi_0\Shape$ of the localized stack with the operator-algebraic set requires knowing
that the condensed shape does not see more (or less) than the operator-algebraic
equivalence, which is a comparison between two homotopy theories that has not been
carried out. In low dimension, where \cite{ogata-classification-review} gives complete
invariants, the conjecture is closest to reach; in higher dimension it is genuinely
open, in step with the general state of the classification problem.

\begin{remark}[The three conjectures as the gap-theoretic Master Conjecture]
\label{rem:master}
Conjectures III-1, III-2, III-3 are the gap-layer components of the program's Master
Conjecture (Part VI \cite{paperSynthesis}). III-1 is the topology of the substack,
III-2 its behavior under profinite disorder, III-3 the comparison of its localized
homotopy type with the established operator-algebraic classification. None can be
proved outright today; each is anchored to a theorem of this paper (III-B, III-A +
\Cref{prop:union}, and \Cref{prop:flow} respectively) that establishes its local or
one-directional content.
\end{remark}

\section[Numerical illustration: the transverse-field Ising chain]{Numerical illustration: the\\ transverse-field Ising chain}
\label{sec:numerics}

We close the mathematical development with a concrete computation that makes the
abstract picture visible: the spectral gap of the transverse-field Ising chain,
computed by exact diagonalization. It renders the gapless discriminant $\Sig$ as a
one-dimensional picture and exhibits the stability that Theorem III-A quantifies.
The accompanying Haskell program (\Cref{sec:numerics-code}) performs the computation
and checks the properties as executable QuickCheck tests.

\subsection{The model}
\label{sec:numerics-model}

On a chain of $N$ spins with open boundary conditions, consider
\begin{equation}
\label{eq:tfim}
H(g)\;=\;-\sum_{i=1}^{N-1} Z_i Z_{i+1}\;-\;g\sum_{i=1}^{N} X_i,
\qquad g\in[0,2],
\end{equation}
where $X_i,Z_i$ are Pauli operators at site $i$. The model has a global $\Z_2$
symmetry $P=\prod_i X_i$, $[H(g),P]=0$, and a quantum phase transition at $g=1$
separating a ferromagnetic phase ($g<1$, two-fold degenerate ground sector in the
thermodynamic limit, the symmetry-broken states) from a paramagnetic phase ($g>1$,
unique ground state). By the Jordan--Wigner transformation the model is a free-fermion
system with single-particle dispersion
\begin{equation}
\label{eq:disp}
\varepsilon_k(g)\;=\;2\sqrt{\,1+g^2-2g\cos k\,},
\end{equation}
whose minimum over $k$ is $\min_k\varepsilon_k(g)=2\lvert 1-g\rvert$, the
thermodynamic single-particle gap. We use open boundary conditions: they keep the
sparse structure and the $\Z_2$ symmetry of the matrix simple and avoid the
fermion-parity (Neveu--Schwarz / Ramond) sector bookkeeping that periodic boundaries
introduce through the Jordan--Wigner transformation. Periodic boundaries would reduce
finite-size edge effects, but the qualitative gap picture---the discriminant at $g=1$
and stability away from it---is independent of the boundary condition in the
thermodynamic limit. Thus
\begin{equation}
\label{eq:sigma}
\Sig=\set{g\,:\,\gapof{H(g)}=0}=\set{g=1},
\end{equation}
a single point of the parameter interval: the gapless discriminant of this
one-parameter family. Crossing $g=1$ is the topological/quantum phase transition;
away from it the system is gapped, with gap growing linearly, $\approx 2\lvert 1-g\rvert$.
The exact gap also exhibits the Kramers--Wannier self-duality
$\min_k\varepsilon_k(g)=g\,\min_k\varepsilon_k(1/g)$, immediate from \eqref{eq:disp}
since $2\lvert 1-g\rvert=g\cdot 2\lvert 1-1/g\rvert$; the Haskell suite of
\Cref{sec:numerics-code} checks this identity.

\subsection{Exact diagonalization and the discriminant}
\label{sec:numerics-ed}

In the computational ($Z$-) basis, \eqref{eq:tfim} is a real symmetric $2^N\times 2^N$
matrix: the $ZZ$ term is diagonal, the transverse field $X_i$ is off-diagonal (a
single spin flip). We diagonalize it exactly for small $N$ by a cyclic Jacobi
eigenvalue routine, sweeping $g$ across $[0,2]$ and recording the low-lying spectrum.
The correctness of the eigensolver is not taken on faith: the QuickCheck properties of
\Cref{sec:numerics-code} check the exact trace identities
\begin{equation}
\label{eq:traces}
\sum_j\lambda_j(g)=\Tr H(g)=0,
\qquad
\sum_j\lambda_j(g)^2=\Tr H(g)^2=\big((N-1)+N g^2\big)\,2^N,
\end{equation}
which follow from $\Tr(Z_iZ_{i+1})=\Tr(X_i)=0$ and the orthogonality of the Pauli
strings, and which a wrong spectrum would violate. The cyclic Jacobi routine drives
the off-diagonal Frobenius norm below $10^{-11}$, and for the dimensions used
($N\le 8$, matrices up to $256\times256$) the eigenvalues are accurate to near machine
precision---the residuals in the identities \eqref{eq:traces} stay below $10^{-9}$
across this range, reaching $\sim 10^{-13}$ at $N=4$. At substantially larger $N$ both the accumulated floating-point error and
the $O(2^{3N})$ cost of dense diagonalization grow, which is why we cap the dense
computation at $N=8$ and read the large-$N$ trend from the analytic reference
\eqref{eq:disp} instead.

The finite-$N$ gap $\gapfin{N}(g)=\lambda_1(g)-\lambda_0(g)$ and the $2$-gap
$\lambda_2(g)-\lambda_0(g)$ together render \eqref{eq:sigma}. In the paramagnetic
phase $g>1$ the ground state is unique and $\gapfin{N}(g)$ is the physical gap,
converging to $2(g-1)$ as $N$ grows. In the ferromagnetic phase $g<1$ the two lowest
levels form the near-degenerate ground sector (their splitting decays exponentially in
$N$), so the physical gap is the $2$-gap---precisely the $m$-gap of
\Cref{rem:why-n} with $m=2$. In both phases the relevant gap develops a sharp minimum
near $g=1$ that deepens with $N$: the finite-size fingerprint of the discriminant point
\eqref{eq:sigma}. \Cref{fig:gap} sketches the shape.

\begin{figure}[h]
\centering
\begin{tikzpicture}[scale=1.05]
  % axes
  \draw[->] (0,0) -- (6.4,0) node[right] {$g$};
  \draw[->] (0,0) -- (0,3.3) node[above] {gap};
  \foreach \x/\l in {0/0, 3/1, 6/2} \draw (\x,0.05) -- (\x,-0.05) node[below] {\l};
  % V-shaped exact gap 2|1-g|, g in [0,2] mapped to x in [0,6]
  \draw[very thick, blue!60!black] (0,3) -- (3,0.06) -- (6,3);
  \node[blue!60!black] at (1.15,2.35) {$2\lvert 1-g\rvert$};
  % finite-N rounded curves near the minimum
  \draw[thick, red!70!black, dashed] plot[smooth] coordinates
    {(0,3.0) (1.5,1.55) (2.4,0.75) (3,0.55) (3.6,0.75) (4.5,1.55) (6,3.0)};
  \node[red!70!black] at (4.9,1.25) {finite $N$};
  % discriminant marker
  \draw[dotted] (3,0) -- (3,3.1);
  \node[below] at (3,-0.35) {\small $\Sig=\{g=1\}$};
  \filldraw[black] (3,0.06) circle (1.4pt);
\end{tikzpicture}
\caption{Schematic of the transverse-field Ising gap across $g\in[0,2]$. The exact
thermodynamic gap (solid) is the ``V'' $2\lvert 1-g\rvert$ of \eqref{eq:disp}, vanishing
only at the discriminant $\Sig=\{g=1\}$ of \eqref{eq:sigma}. Finite-$N$ gaps (dashed)
round the minimum and deepen toward the ``V'' as $N$ grows. Away from $\Sig$ the gap is
bounded below and the family is uniformly gapped on any closed sub-interval avoiding
$g=1$.}
\label{fig:gap}
\end{figure}

\subsection{Perturbative stability away from criticality}
\label{sec:numerics-pert}

To illustrate Theorem III-A and the stability picture, the program performs a
perturbation experiment. Fix $g=1.5$, deep in the gapped paramagnetic phase, and add a
random local field $V=\lambda\sum_i s_i Z_i$ with $s_i\in\{\pm 1\}$ and small
$\lambda$. Exact diagonalization confirms
\begin{equation}
\label{eq:lip-numeric}
\big\lvert\gapfin{N}(H(1.5)+V)-\gapfin{N}(H(1.5))\big\rvert\;\le\;2\,\norm{V}
\;\le\;2\lambda N,
\end{equation}
the finite-volume Lipschitz bound of Theorem III-A (here $C_\Lambda$ absorbed into the
per-site normalization). The gap stays bounded away from $0$ for all sampled
perturbations: away from the discriminant the gapped locus is open and stability holds,
exactly as Theorem III-A guarantees at finite volume. Near $g=1$, by contrast, the same
size of perturbation moves the small gap by a comparable amount and can push the finite
system across the rounded minimum---the numerical trace of why openness of the
substack is a statement scoped away from the discriminant.

\subsection{What the finite computation cannot do}
\label{sec:numerics-cannot}

It is essential to state what \Cref{fig:gap} does \emph{not} establish, on pain of
contradicting \Cref{sec:wall}. No finite-$N$ computation decides the thermodynamic gap.
For the transverse-field Ising chain we happen to know the answer analytically---the
model is exactly solvable, \eqref{eq:disp}---so the discriminant \eqref{eq:sigma} is
certain; but that certainty comes from the exact Jordan--Wigner solution, not from the
diagonalization. For a general translation-invariant family the analogous finite-$N$
data would be genuinely uninformative about the thermodynamic limit, because by
\Cref{thm:cpw} no finite computation, however large $N$, can decide whether the gap
survives. The Ising chain is an illustration of the \emph{structure} of the gapped
substack---its discriminant, its stability away from criticality, the Lipschitz bound---%
not a method for detecting gaps. The distinction is the entire moral of Part III\@.

\subsection{The Haskell computation}
\label{sec:numerics-code}

The computation is implemented in Haskell in \texttt{src/spectral-gap-stability/}. The
module \texttt{TFIM} builds the Hamiltonian \eqref{eq:tfim} as a real symmetric matrix
and diagonalizes it by cyclic Jacobi rotations; \texttt{Main} sweeps $g$ over $[0,2]$
for $N\in\{4,6,8\}$, prints the gap table, runs the perturbation experiment of
\Cref{sec:numerics-pert}, and exits with status $0$ on success. The module
\texttt{Properties} encodes as QuickCheck properties: the exact trace identities
\eqref{eq:traces} (which double as a validation of the eigensolver); symmetry of the
Hamiltonian matrix; positivity of the gap in the paramagnetic phase at fixed $N$
(\Cref{sec:numerics-ed}); the finite-volume Lipschitz bound \eqref{eq:lip-numeric} of
Theorem III-A; and the monotone growth of the gap with $g$ in the paramagnetic phase.
The properties are stated one per claim and cite the corresponding equation or theorem,
so that the executable tests track the paper's assertions.

\section{Discussion}
\label{sec:discussion}

\subsection{What the condensed paradigm contributes here}
\label{sec:disc-contribute}

Part III does not prove a new stability theorem; the analytic content is
\cite{bhm-stability,bravyi-hastings-shortproof,michalakis-zwolak,nsy-bulkgap,hastings-koma,nachtergaele-sims-clustering}.
Nor is the idea of studying the moduli space of gapped Hamiltonians and its topology
new: it is developed, in ordinary topology and with effective-field-theory methods, by
Hsin, Wang and collaborators \cite{hsin-wang-moduli}. What the condensed framing
contributes is organizational, and it is not empty. First,
it makes ``uniformly gapped'' a property of a \emph{family} in a category where
families, disorder hulls, and inverse limits are first-class objects: the $\einf_s$ of
\Cref{def:gap} is the condition that a subobject of a condensed stack is cut out
correctly, and Theorem III-C shows this is the condition with structural teeth.
Second, it identifies exactly which classical theorems give \emph{openness} of the
distinguished subobject (Theorem III-B) and on which stratum, so that the downstream
constructions of Parts IV and V know precisely what they may assume. Third, it places
the undecidability wall where it belongs (as a statement about global decidability
that leaves the local, stratified, functorial theory intact) rather than letting it
either be ignored or overstated into paralysis.

\subsection{The two walls}
\label{sec:disc-walls}

The program has two hard external constraints. The first is the undecidability of the
gap \cite{cpw-undecidable}, which is the governing constraint of this paper: it is why
existence of a thermodynamic gap is a hypothesis and not a theorem, why openness is a
stratum statement, and why \Cref{sec:numerics-cannot} insists that no finite
computation detects the gap. The second, the Kapustin--Fidkowski obstruction \cite{kapustin-fidkowski} to
commuting-projector realizations of chiral phases, governs Part V
\cite{paperRealizability} and does not bear directly here, except to note that the
frustration-free / LTQO stratum of \Cref{def:stratum} (on which we can prove
stability) is close to the commuting-projector world where the second wall bites.
Chiral phases (nonzero Hall conductance, \cite{kapustin-sopenko-noether}) are gapped
but are not frustration-free commuting-projector systems, so they lie outside
$\mathfrak{S}$; our stability theorem says nothing about them, and it should not,
because their stability is a different and harder question.

\subsection{The profinite-disorder thread}
\label{sec:disc-disorder}

The disorder hull $\dis=Q^{\Z^d}$ is where the condensed and profinite structure is
not a repackaging of a smooth manifold but a match to the physical configuration space
\cite{bellissard-ncg-qhe,kellendonk-tilings,prodan-sb}. Conjecture III-2 is the
gap-layer instance of the program's recurring descent question: does a
finite-resolution property, holding compatibly at every finite quotient of $\dis$ with
uniform constants, descend to the profinite limit? For gaps this is delicate precisely
because the gap is thermodynamic. That the question is even well posed (that a
disordered family is literally a $\dis$-point of $\Ham_{d,G}$, and uniform gappedness a
condition on that point) is a contribution of the condensed viewpoint, and settling it
would connect the gap-labelling of \cite{bellissard-ncg-qhe} to the substack
$\Gap_{d,G}$ directly.

\subsection{Limitations}
\label{sec:disc-limits}

We list the limitations plainly. Theorem III-B is conditional on membership in a
frustration-free / LTQO stratum, which is undecidable to verify in general and is a
proper subclass of gapped systems. Theorem III-C assumes the uniform Lieb--Robinson
data of Part I, which restricts to the $\Ffun$-function locality class and excludes
genuinely long-range interactions. \Cref{prop:lsc} gives only lower semicontinuity, and
deliberately so. The three conjectures are unproven, and Conjecture III-1 in particular
cannot be strengthened to openness in the full $\Ham_{d,G}$ without contradicting
\cite{cpw-undecidable}. None of these limitations is incidental; each marks a boundary
that the undecidability of the gap, or the restriction to a provable stability class,
places on what can be claimed.

\section{Conclusion}
\label{sec:conclusion}

The uniformly gapped substack $\Gap_{d,G}$ is the object on which the condensed theory
of topological phases is built, and it is governed by a hard theorem: the spectral gap
is undecidable, so membership cannot be the output of a computation and must be taken
as a defining hypothesis. Within that constraint we proved what can be proved. At
finite volume the gap is Lipschitz and the gapped locus open (Theorem III-A)\@. On the
frustration-free / LTQO stratum the classical stability theorems give a
size-independent perturbation threshold, which we recast as condensed-openness of
$\Gap_{d,G}$ along quasi-local directions (Theorem III-B), with an explicit spectral
flow realizing the equivalence class $\Weq$ (\Cref{prop:flow}). A uniform gap forces a
uniform correlation length (Theorem III-C), the structural property that distinguishes
uniformly gapped families from pointwise-gapped ones and that makes $\einf_s$ the
correct condition in the definition of the substack. Beyond the stratum we recorded
three conjectures: openness on the physical stratum with finite-quotient detection
(III-1), uniform-gap descent along profinite disorder hulls (III-2), and the
identification of quasi-adiabatic path components with the operator-algebraic gapped
phases (III-3), each anchored to a theorem of this paper and each scoped to respect the
undecidability wall. The transverse-field Ising computation made the discriminant
$\Sig=\{g=1\}$ and the stability away from it visible, while insisting that no finite
computation detects the thermodynamic gap.

The substack is now a well-understood object on its provable stratum and a precisely
delineated open problem beyond it. Part IV takes $\Gap_{d,G}$ as given and
group-completes its invertible sector into the condensed phase spectrum
$\IPcond_{d,G}$; Part V asks which classes that spectrum records are realized by
uniformly gapped lattice families. Both rest on the stability established here, and
both inherit its scope.

\section*{Code availability}
\addcontentsline{toc}{section}{Code availability}
The Haskell source for the transverse-field Ising computation of \Cref{sec:numerics}
(exact diagonalization, the base-only cyclic Jacobi eigensolver, and the QuickCheck
property suite) is available at
\href{https://github.com/YonedaAI/topological-phases-of-matter}{github.com/YonedaAI/topological-phases-of-matter},
in \texttt{src/spectral-gap-stability/}.

\bibliographystyle{plain}
\bibliography{references}

\end{document}
